
\documentclass{article}
\usepackage{colortbl}
\usepackage{makecell}
\usepackage{multirow}
\usepackage{supertabular}

\begin{document}

\newcounter{utterance}

\centering \large Interaction Transcript for game `cladder', experiment `full\_v1.5\_default', episode 5136 with Qwen3.6{-}27B.
\vspace{24pt}

{ \footnotesize  \setcounter{utterance}{1}
\setlength{\tabcolsep}{0pt}
\begin{supertabular}{c@{$\;$}|p{.15\linewidth}@{}p{.15\linewidth}p{.15\linewidth}p{.15\linewidth}p{.15\linewidth}p{.15\linewidth}}
    \# & \multicolumn{2}{c}{Player} && \multicolumn{2}{c}{Game Master} \\
    \hline

    \theutterance \stepcounter{utterance}  
    & & & \multicolumn{4}{p{0.6\linewidth}}{
        \cellcolor[rgb]{0.9,0.9,0.9}{
            \makecell[{{p{\linewidth}}}]{
                \texttt{\tiny{[P1$\langle$GM]}}
                \texttt{You are an expert in causal inference. The following question is not a typical commonsense query, but rather a meticulously designed question created by a professor specializing in causal inference, intended to assess the students' mastery of the course content.} \\
\texttt{Imagine a self{-}contained, hypothetical world with only the following conditions, and without any unmentioned factors or causal relationships: Unobserved confounders has a direct effect on education level and foot size. Proximity to a college has a direct effect on education level. Education level has a direct effect on foot size. Unobserved confounders is unobserved. The overall probability of college degree or higher is 47\%. The probability of high school degree or lower and large feet is 32\%. The probability of college degree or higher and large feet is 10\%. Is the chance of large feet smaller when observing college degree or higher?} \\
\texttt{Start your answer with 'Yes' or 'No', followed by additional reasoning or evidence to support your explanation.} \\
            }
        }
    }
     \\ \\

    \theutterance \stepcounter{utterance}  
    & \multicolumn{4}{p{0.6\linewidth}}{
        \cellcolor[rgb]{0.561,0.737,0.561}{
            \makecell[{{p{\linewidth}}}]{
                \texttt{\tiny{[P1$\rangle$GM]}}
                \texttt{Yes} \\
\\ 
\texttt{**Reasoning:**} \\
\\ 
\texttt{To determine if the chance of large feet is smaller when observing a college degree or higher, we need to compare the conditional probability of having large feet given a college degree or higher, \$P(\textbackslash{}text\{Large Feet\} \textbackslash{}mid \textbackslash{}text\{College Degree+\})\$, with the conditional probability of having large feet given a high school degree or lower, \$P(\textbackslash{}text\{Large Feet\} \textbackslash{}mid \textbackslash{}text\{High School{-}\})\$. Alternatively, we can simply calculate the specific probability requested and compare it to the overall prevalence or the other group, but the question specifically asks if the chance is *smaller* in the college group. Usually, this implies a comparison to the alternative group (High School or lower) or the baseline. Let's calculate the specific probabilities.} \\
\\ 
\texttt{**1. Define the Variables and Given Probabilities:**} \\
\texttt{*   Let \$E\$ be the event of having a College Degree or higher.} \\
\texttt{*   Let \$\textbackslash{}neg E\$ be the event of having a High School Degree or lower.} \\
\texttt{*   Let \$F\$ be the event of having Large Feet.} \\
\texttt{*   Let \$\textbackslash{}neg F\$ be the event of not having Large Feet.} \\
\\ 
\texttt{Given:} \\
\texttt{*   \$P(E) = 0.47\$} \\
\texttt{*   Therefore, \$P(\textbackslash{}neg E) = 1 {-} 0.47 = 0.53\$} \\
\texttt{*   \$P(\textbackslash{}neg E \textbackslash{}cap F)} \\
            }
        }
    }
    & & \\ \\

    \theutterance \stepcounter{utterance}  
    & & & \multicolumn{2}{p{0.3\linewidth}}{
        \cellcolor[rgb]{0.9,0.9,0.9}{
            \makecell[{{p{\linewidth}}}]{
                \texttt{\tiny{[GM$|$GM]}}
                \texttt{Yes} \\
            }
        }
    }
    & & \\ \\

    \theutterance \stepcounter{utterance}  
    & & & \multicolumn{2}{p{0.3\linewidth}}{
        \cellcolor[rgb]{0.9,0.9,0.9}{
            \makecell[{{p{\linewidth}}}]{
                \texttt{\tiny{[GM$|$GM]}}
                \texttt{yes} \\
            }
        }
    }
    & & \\ \\

    \theutterance \stepcounter{utterance}  
    & & & \multicolumn{2}{p{0.3\linewidth}}{
        \cellcolor[rgb]{0.9,0.9,0.9}{
            \makecell[{{p{\linewidth}}}]{
                \texttt{\tiny{[GM$|$GM]}}
                \texttt{game\_result = WIN} \\
            }
        }
    }
    & & \\ \\

\end{supertabular}
}

\end{document}
